% !PRINTED 714
\weberpage{714}
Wenn aber $\lambda=1$ ist, so muss $\beta=0$ sein, und es sind
noch zwei Fälle möglich:
\[
\begin{aligned}
&\text{entweder}\quad f=1,\ \alpha=0,\ p\equiv\phantom{-}1\pmod{m},\ \lambda=f\\
&\text{oder}\qquad\quad\, f=2,\ \alpha=1,\ p\equiv -1\pmod{m},\ \lambda=\tfrac12 f.
\end{aligned}
\]
Nur in diesem letzten Falle gehört $p$ zur ersten Art, und es
folgt also, dass auch hier allgemein $\Theta^{\beta}$ eine
primitive $f_1^{\text{te}}$ Einheitswurzel ist. Hiermit ist dann
der Satz bewiesen.

Dieser Satz gestattet uns die Anwendung der Formel \S.\,193, (9),
die hier ergiebt:
\weq{8}{(1-p^{-sf_1})^{e_1} = \prod^{\chi_1}\bigl[1-\chi_1(p)p^{-s}\bigr],}
worin sich das Product auf alle Charaktere $\chi_1$ der Gruppe
$C_1$ erstreckt, und nun lässt sich die Function $\Phi_1(s)$
ebenso wie $\Phi(s)$ im \S.\,193 weiter entwickeln. Man erhält auf
demselben Wege, wie dort, die Formel (11):
\weq{9}{\Phi_1(s) = \frac{1}{1-q^{-s}}\;\prod^{\chi_1}\sum^{n}\frac{\chi_1(n)}{n^s},}
und wenn man zur Grenze $s=1$ übergeht,
\weq{10}{g_1 h_1 = \prod^{\chi_1}\sum^{n}\frac{\chi_1(n)}{n},}
worin aber jetzt das Product auf alle Charaktere der Gruppe $C_1$
\gesperrt{mit Ausnahme des Hauptcharakters} zu erstrecken ist. Dieses Product
ist also ein Theil des Productes, durch welches im \S.\,193 die
Zahl $gh$ ausgedrückt ist.

Die Zahlen $g$, $g_1$ ergeben sich aus dem im \S.\,190, (11)
angegebenen allgemeinen Ausdrucke:
\weq{11}{\frac{2^{\nu}\pi^{n-\nu}L}{w\sqrt{\pm\Delta}}.}
Da, wie im \S.\,178 nachgewiesen ist, die Einheiten in $\Omega_m$
und $H_m$, von den Einheitswurzeln abgesehen, dieselben sind, so
ist auch ein Fundamentalsystem von Einheiten in $H_m$ zugleich ein
Fundamentalsystem in $\Omega_m$.

Bedeutet $\varepsilon_1,\varepsilon_2,\dots,\varepsilon_{\nu-1}$
daher ein Fundamentalsystem reeller Einheiten, und sind
$\varepsilon_{i,1},\varepsilon_{i,2},\dots,\varepsilon_{i,\nu}$ die
conjugirten Einheiten zu $\varepsilon_i$, so sind die conjugirten
Logarithmen im Körper $H_m$:
\[ \log|\varepsilon_{i,1}|,\ \log|\varepsilon_{i,2}|,\ \dots,\ \log|\varepsilon_{i,\nu}|, \]
dagegen im Körper $\Omega_m$, in dem nur conjugirt imaginäre Paare
vorkommen (\S.\,185),
\[ 2\log|\varepsilon_{i,1}|,\ 2\log|\varepsilon_{i,2}|,\ \dots,\ 2\log|\varepsilon_{i,\nu}|. \]
